\section{The charge carrier and its slots}

\begin{proposition}[Charge]\label{thm:charge}\lean{Infomaterial.Charge}\leanok\uses{thm:floor, thm:corner}
A \emph{charge} is the four-point value object of positing. Formally it is
the same four-corner object as \texttt{Articulating.Corner}, now read as
value rather than articulation status.
\end{proposition}

\begin{proposition}[Evidence slots]\label{thm:charge-slots}\lean{Infomaterial.Charge.pos}\leanok\uses{thm:charge, thm:status-correspondence}
Each charge has a positive and a negative Boolean slot. The maps
\texttt{pos}, \texttt{neg}, and \texttt{ofSlots} identify the four charges
with the two-slot object $\{0,1\}^2$.
\end{proposition}

\begin{proposition}[The four distinguished charges]\label{thm:four-charges}\lean{Infomaterial.Charge.kBot}\leanok\uses{thm:charge-slots}
The four slots receive their bilattice names:
$\bot_k$, $\top_k$, $t$, and $f$. They are the same four points read by the
knowledge and truth orders.
\end{proposition}
